\section{Fichiers NetCDF}

\newcommand{\prodcdf}{\textbf{\texttt{prod\_CDF.py}} }
\newcommand{\netcdf}{\texttt{NetCDF} }
\newcommand\myfontsize{\fontsize{9pt}{10pt}\selectfont}

%%%\subsection{\prodcdf}
Le programme \prodcdf convertit un ensemble de
fichiers-vecteurs \textbf{synchronisés sur une même base de temps} d'une même
plateforme en un (ou plusieurs) fichier(s) au format \netcdf.\\

\subsection{Lancement}
\prodcdf doit être lancer en ligne de commande dans un terminal. \\
Pour obtenir de l'aide sur \prodcdf, lancer la commande :
\fbox{\texttt{prod\_CDF.py -h}}

\medbreak
\textbf{\textit{Note : le programme lancé sous \texttt{aito}, onglet
\texttt{Exportation} est : \texttt{prod\_CDF\_histo.py}. \\
Lancer la commande \fbox{\texttt{prod\_CDF\_histo.py -h}} pour obtenir de
l'aide.}
}


\subsection{Manuel d'utilisation}
Par convention, l'extension des fichiers générés est : \texttt{.nc}.

\medbreak
Il utilise plusieurs sources d'information :
\begin{itemize}
\item
  la ligne d'appel
\item
  le fichier de configuration par défaut
\item
  le fichier de configuration
\item
  le fichier AIDA \texttt{plate-formes}
\item
  le fichier de sélection de vecteur
\item
  les fichiers Vecteur AIDA
\item
  les fichiers .INF
\end{itemize}

\medbreak
Il s'appuie sur plusieurs conventions et documentations :
\begin{itemize}
\item
  \texttt{Normalisation\_des\_métadonnées\_NetCDF\_TRAMM\_usage.pdf}
  (la documentation TRAMM)
\item
  \url{http://cfconventions.org/Data/cf-standard-names/current/build/cf-standard-name-table.html}
\item
  \url{https://wiki.esipfed.org/Attribute\_Convention\_for\_Data\_Discovery\_1-3 }
\item
  \url{https://cfconventions.org/Data/cf-conventions/cf-conventions-1.10/cf-conventions.html}
\item
  \url{https://docs.unidata.ucar.edu/udunits/current/}
\end{itemize}


\subsection{Syntaxe de la ligne d'appel}

La syntaxe d'appel de \prodcdf est la suivante :
\medbreak

\fbox{ \myfontsize\texttt{\prodcdf [-h] [-r dir] [-S select\_file] [-v version] -c config\_file prod\_deb prod\_fin} }
\medbreak

avec :
\medbreak

\begin{center}
\tablefirsthead{}
\tablehead{\hline
  \textbf{option} &
  \textbf{défaut} &
  \textbf{commentaire} \\}
\tabletail{\hline}
\tablelasttail{}
  \begin{supertabular}{|p{4.5cm}|p{1.2cm}|p{9cm}|}
\hline
\textbf{option} &
\textbf{défaut} &
\textbf{commentaire} \\
\hline

\texttt{-h} &
&
affichage de l'aide.
\\ \hline

\texttt{-r dir} \smallbreak
\texttt{-{}-rep\_out dir} \smallbreak
\texttt{-o dir} \smallbreak
\texttt{-{}-out\_dir dir} &
\texttt{./} &
nom du repertoire de sortie
\\ \hline

\texttt{-S select\_file} \smallbreak
\texttt{-{}-select select\_file} &
&
chemin du fichier de sélection des vecteurs \smallbreak
pour les paramètres. \smallbreak
voir \nameref{ref:vector_select}
\\ \hline

\texttt{-c config\_file} \smallbreak
\texttt{-{}-config config\_file} &
&
chemin du fichier de configuration.
\\ \hline

\texttt{-v version} \smallbreak
\texttt{-{}-version version} &
\texttt{1} &
numéro de version des fichiers à produire
\\ \hline

prod\_deb \smallbreak
prod\_fin &
&
sous la forme JJMMAAAA : date de début et de fin à traiter. \smallbreak
sous la forme acqid    : premier et dernier acqid à traiter.
\\ \hline

\end{supertabular}
\end{center}

\begin{warning}
  \begin{itemize}
  \item
    La version historique \texttt{prod\_CDF\_histo.py} est lancée si le
    fichier de configuration ne contient pas de section
    \texttt{[DIRECTIVE]} (voir \nameref{ref:configuration}).
  \item
    L'option \texttt{-S\ select} est incompatible avec la configuration
    \texttt{mode\ =\ periode} (voir \nameref{ref:configuration})
  \item
    \textbf{Le numéro de version ne peut être précisé que par cette
    option}.\\
    Il renseigne l'attribut global \texttt{product\_version} du fichier
    \netcdf.\\
    L'option \texttt{product\_version} dans le fichier de configuration
    est ignorée.
  \item
    Le nom du fichier \netcdf généré dépend de différentes options
    fournies. Voir \nameref{ref:nom_NetCDF}.
  \end{itemize}
\end{warning}

\subsection{Configuration par défaut}

Les informations à inscrire dans les fichiers \netcdf qui
ne changent que très rarement pour une équipe peuvent être définies dans
un fichier de configuration par défaut.
\medbreak

Son nom et sa localisation sont :
\texttt{\textbf{\$AIDA\_HOME/profils/\$AIDA\_PROFIL/CDF.defaults}}
\medbreak

Ce fichier est de type INI (voir \nameref{sec:anx_ini}). \\
\medbreak

Il ne contient qu'une seule section : \texttt{[GENERAL]}
\medbreak


\begin{warning}
Les valeurs des options inscrites dans le fichier de configuration écrassent
celles inscrites dans le fichier de configuration par défaut.\\
Ex :

\begin{longtable}[]{@{}lll@{}}
\toprule
\begin{minipage}[b]{0.30\columnwidth}\raggedright
configuration par défaut\strut
\end{minipage} & \begin{minipage}[b]{0.30\columnwidth}\raggedright
configuration\strut
\end{minipage} & \begin{minipage}[b]{0.30\columnwidth}\raggedright
\netcdf\strut
\end{minipage}\tabularnewline
\midrule
\endhead
\begin{minipage}[t]{0.30\columnwidth}\raggedright
processing\_level = L2\strut
\end{minipage} & \begin{minipage}[t]{0.30\columnwidth}\raggedright
processing\_level = L1\strut
\end{minipage} & \begin{minipage}[t]{0.30\columnwidth}\raggedright
:processing\_level = "L1" ;\strut
\end{minipage}\tabularnewline
\bottomrule
\end{longtable}
\end{warning}

\subsubsection{Remarque}

Pour plus de lisibilité du fichier de configuration, les longues valeurs
des options peuvent être enregistrées sur plusieurs lignes (le caractère
de retour à la ligne (\texttt{\textbackslash{}n}) sera inscrit dans le
fichier \netcdf).

Ex :\\
Si le fichier de configuration contient :

\begin{verbatim}
contributor_name = prénom1: nom1: prénom1.nom1@x.fr: institut1;
                   prénom2: nom2: prénom2.nom2@x.fr: institut2;
\end{verbatim}

le fichier \netcdf contiendra :

\begin{verbatim}
string :contributor_name = "prénom1: nom1: prénom1.nom1@x.fr: institut;\nprénom2: [...]
\end{verbatim}

\subsection{Configuration}\label{ref:configuration}

Les informations spécifiques à la génération d'un fichier
\netcdf sont définies dans un fichier de configuration.
\medbreak

L'extension de son nom est préférentiellement \texttt{.ccdf}.
\medbreak

Ce fichier est de type INI (voir \nameref{sec:anx_ini}). \\
\medbreak

Ce fichier est constitué de plusieurs sections :
\medbreak

\begin{itemize}
\item
  une section \texttt{[DIRECTIVES]}\\
  elle rassemble des informations permettant à \prodcdf de
  travailler (en particulier des informations utiles pour générer le nom
  des fichiers \netcdf) ou pour générer la partie
  \textit{dimensions} du fichier \netcdf.
\item
  une section \texttt{[GENERAL]}\\
  on \textbf{n'y trouve que} des informations qui seront inscrites sous
  le même nom dans le fichier \netcdf dans la partie
  \textit{global attributes}.
\item
  une section par paramètre\\
  pour chaque paramètre, on trouve des informations qui seront inscrites
  dans la partie \textit{variables} du fichier \netcdf ou
  dans la partie \textit{dimensions} .
\end{itemize}
\medbreak

Hors la section \texttt{[DIRECTIVES]}, les sections peuvent
également contenir des \textit{options libres} (voir
\nameref{ref:options_libres})

\subsubsection{Section \texttt{[DIRECTIVES]}}

Cette section des \textbf{obligatoire}. \\
Si elle est absente, \prodcdf lancera la version historique :
\texttt{prod\_CDF\_histo.py}\strut

\begin{center}
\tablefirsthead{}
\tablehead{\hline
  \textbf{option} &
  \textbf{nature}

  \textbf{défaut} &
  \textbf{commentaire} &
  \textbf{exemple} \\}
\tabletail{\hline}
\tablelasttail{}
\begin{supertabular}{|p{3.85cm}|p{1.7cm}|p{4.6cm}|p{5.3cm}|}
\hline
\textbf{option} &
\textbf{nature}

\textbf{défaut} &
\textbf{commentaire} &
\textbf{exemple} \\
\hline

prefix &
obligatoire &
utilisé dans la génération du nom du fichier \netcdf \smallbreak
voir \nameref{ref:nom_NetCDF} &
\textit{CNRM\_MTO\_L2A-30min-Mean}\\
\hline

plateforme &
obligatoire &
\textbf{il ne s'agit pas du champ "platform" du fichier \netcdf} mais de la
plateforme AIDA &
\textit{toulouse-flux} \\
\hline

reftime &
obligatoire &
référence temporelle du fichier \netcdf  \smallbreak
\textbf{et des vecteurs à traiter} &
\textit{rft\_time\_REG\_\%1800} \\
\hline

mode &
obligatoire &
periode ou jour ou mobile \smallbreak
voir \nameref{ref:mode} &
\textit{mobile} pour les plateformes mobiles \\
\hline

date\_format &
facultatif \smallbreak
\%Y\%m\%d &
utilisé dans la génération du nom du fichier \netcdf &
\textit{\%Y-\%m-\%d} \\
\hline

replace\_missing\_vector &
facultatif \smallbreak
False &
les données absentes pour une acquisition sont remplacées par Nan &
\textit{True} \\
\hline

datation\_convention &
facultatif \smallbreak
fin &
utilisé pour la création des bornes temporelles &
\textit{debut} \\
\hline

dimensions &
facultatif \smallbreak
Unlimited &
Unlimited ou Limited

champ \texttt{dimensions} du fichier \netcdf &
\textit{Limited} \\
\hline

\end{supertabular}
\end{center}

\begin{warning}
L'absence d'une option obligatoire de la section
\texttt{[DIRECTIVES]} génère l'affichage d'une erreur (\texttt{+++})
et l'arrêt du programme.
\end{warning}

\begin{warning}
Si la référence temporelle de l'un des vecteurs à traiter n'est pas cohérente
avec le vecteur nommé par l'option \texttt{reftime}, la génération
affiche un message d'erreur (\texttt{+++}) et s'arrête.
\end{warning}


\paragraph{Nom du fichier \netcdf} : \label{ref:nom_NetCDF}

Le nom du fichier \netcdf généré dépend :

\begin{itemize}
\item
  des options fournies dans le fichier de configuration par défaut
\item
  des options fournies dans le fichier de configuration
\item
  de la version précisée par l'option -v de la ligne d'appel
\item
  de l'acqid (pour le champ \{acqid\_date\} du mode \texttt{jour})
\item
  de la date de l'acquisition lues dans le fichier .INF (pour le champ
  \{acqid\_date\} du mode \texttt{mobile})
\end{itemize}

\medbreak
Il est construit de la façon suivante, selon les informations positionnées :

\begin{itemize}
\item
  si mode vaut \textbf{jour} : \\
    \{prefix\}\_\{acqid\_date\}\_V\{version\} ou \\
    \{prefix\}\_\{acqid\_date\}\_\{processing\_level\}\_V\{version\}
    \medbreak

\item
  si mode vaut \textbf{periode} : \\
    \{prefix\}\_\{date\_start\}\_\{date\_end\}\_V\{version\} ou \\
    \{prefix\}\_\{date\_start\}\_\{date\_end\}\_\{processing\_level\}\_V\{version\}
    \medbreak

\item
  si mode vaut \textbf{mobile} : \\
    \{prefix\}\_\{acqid\_date\}\_\{acqid\}\_V\{version\} ou \\
    \{prefix\}\_\{acqid\_date\}\_\{acqid\}\_\{processing\_level\}\_V\{version\} \\
\end{itemize}


\paragraph{Option mode} : \label{ref:mode}

L'option \texttt{mode} permet de spécifier le type de fichier
\netcdf à générer.
\medbreak

Cette option peut prendre une des trois valeurs : \texttt{periode},
\texttt{jour}, \texttt{mobile}.

\begin{itemize}
\item
  \texttt{mode = periode} :

    \prodcdf produit un seul fichier \netcdf couvrant la période définie par
    les arguments \texttt{prod\_deb} \texttt{prod\_fin} de la ligne d'appel.

\item
  \texttt{mode = jour} :
    
    \prodcdf produit un fichier \netcdf pour chaque acquisition de la période
    définie par les arguments \texttt{prod\_deb} \texttt{prod\_fin} de la ligne
    d'appel.

\item
  \texttt{mode = mobile} :
    \prodcdf produit un fichier \netcdf de type \textit{Trajectory} pour chaque
    acquisition de la période définie par les arguments \texttt{prod\_deb}
    \texttt{prod\_fin} de la ligne d'appel.
\end{itemize}

\subsubsection{Section \texttt{[GENERAL]}}

Les options \textcolor{orange}{colorées en orange} sont spécifiques au mode
\textcolor{orange}{\texttt{mobile}}.

\begin{center}
\tablefirsthead{}
\tablehead{\hline
  \textbf{option} &
  \textbf{nature}

  \textbf{défaut} &
  \textbf{commentaire} \\}
\tabletail{\hline}
\tablelasttail{}
\begin{supertabular}{|p{4.6cm}|p{4cm}|p{7.5cm}|}
\hline
\textbf{option} &
\textbf{nature}

\textbf{défaut} &
\textbf{commentaire}
\\ \hline

title &
obligatoire &
\\ \hline

summary &
obligatoire &
\\ \hline

Conventions &
obligatoire  \smallbreak
\textit{CF-1.8, ACDD-1.3} &
\\ \hline

keywords &
obligatoire &
devraient être issus de \smallbreak
\url{https://gcmd.earthdata.nasa.gov/KeywordViewer/scheme/all?gtm\_scheme=all}
\\ \hline

keywords\_vocabulary &
obligatoire \smallbreak
\textit{GCMD: }\url{https://gcmd.earthdata.nasa.gov/KeywordViewer/scheme/all?gtm\_scheme=all} &
\\ \hline

data\_policy &
obligatoire \smallbreak
\textit{""} &
il est souhaitable d'indiquer un data\_policy
\\ \hline

processing\_level &
obligatoire \smallbreak
\textit{L2} &
automatique avec fichier .defaults
\\ \hline

source &
facultatif &
cf. documentation TRAMM : \smallbreak
\textit{surface observation, radiosonde, airborne measurements}
\\ \hline

platform &
souhaité si défini sous \smallbreak
\url{https://skosmos.aeris-data.fr/aeris\_pf/en/}
\\ \hline

doi &
obligatoire &
\\ \hline

creator\_name &
obligatoire &
\\ \hline

source\_config &
automatique &
le nom du fichier de configuration
\\ \hline

institution &
souhaitable &
automatique avec fichier .defaults
\\ \hline

creator\_institution &
souhaitable &
automatique avec fichier .defaults
\\ \hline

creator\_email &
souhaitable &
automatique avec fichier .defaults
\\ \hline

contributor\_role &
souhaitable &
cf. documentation TRAMM : \smallbreak
\textit{principalinvestigator resourceprovider contributor}
\textit{owner funder sponsor editor coauthor custodian pointofcontact author}
\textit{originator distributor processor stakeholder}
\textit{collaborator metadatapointofcontact publisher}
\textit{mediator user rightsholder metadatapointofcontact}
\\ \hline

contributor\_name &
souhaitable &
cf. documentation TRAMM : \smallbreak
\textit{prenom:nom:email:institution}
\\ \hline

id &
automatique &
\{prefix\}\_\{processing\_level\}
\\ \hline

project &
souhaitable &
\\ \hline

comment &
souhaitable &
cf. documentation TRAMM
\\ \hline

coordinates &
préférable pour une plateforme mobile &
voir \nameref{ref:geospatial}
\\ \hline

featureType &
automatique \smallbreak
(si l'option `coordinates` est fournie) \textit{Trajectory} &
voir \nameref{ref:featureType}
\\ \hline

history &
souhaitable et \smallbreak
automatique &
\prodcdf ajoute sa ligne de lancement au champ fourni
\\ \hline

date\_created &
automatique &
date et heure d'exécution de \prodcdf
\\ \hline

date\_modified &
automatique &
date et heure d'exécution de \prodcdf
\\ \hline

time\_coverage\_start &
automatique &
première datation
\\ \hline

time\_coverage\_end &
automatique &
dernière datation
\\ \hline

time\_coverage\_resolution &
automatique \smallbreak
PT\{1./freq\_ech\}S &
ex : PT1800.000000S
\\ \hline

geospatial\_lat\_min \smallbreak
geospatial\_lat\_max \smallbreak
geospatial\_lon\_min \smallbreak
geospatial\_lon\_max \smallbreak
\textcolor{orange}{geospatial\_lat\_units} \smallbreak
\textcolor{orange}{geospatial\_lon\_units} \smallbreak
\textcolor{orange}{geospatial\_vertical\_min} \smallbreak
\textcolor{orange}{geospatial\_vertical\_max} \smallbreak
\textcolor{orange}{geospatial\_vertical\_units} \smallbreak
\textcolor{orange}{geospatial\_vertical\_positive} &
automatique &
A partir du fichier plate-formes pour les plateformes fixes. \smallbreak
A partir des vecteurs de l'option \texttt{coordinates} pour les plateformes mobiles. \smallbreak
voir \nameref{ref:geospatial}
\\ \hline

acqid / \textcolor{orange}{flight\_id} &
automatique &
\textcolor{orange}{flight\_id} pour toutes les plateformes mobiles même les voitures...
\\ \hline

\textcolor{orange}{flight\_date} &
automatique &
A partir du fichier \texttt{.inf}.
\\ \hline

\textcolor{orange}{time\_take\_off} &
automatique &
A partir du fichier \texttt{.inf}.
\\ \hline

\textcolor{orange}{time\_landing} &
automatique &
A partir du fichier \texttt{.inf}.
\\ \hline

\end{supertabular}
\end{center}

\begin{warning}
  Les options qui ne sont pas explicitées ci-dessus sont inscrites dans la
  partie \textit{global attributes} du fichier \netcdf.
\end{warning}


\paragraph{Champs \texttt{geospatial}} : \label{ref:geospatial}

S'ils ne sont pas initialisés manuellement dans le fichier de
configuration, la génération des champs \textit{geospatial\_XXX} dans le
fichier \netcdf se fait de la façon suivante :

\begin{itemize}
\item
  pour une plateforme fixe :\\
  les champs \textit{geospatial\_XXX} sont initialisés avec les
  informations trouvées dans le fichier
  \texttt{\$AIDA\_HOME/profils/\$AIDA\_PROFIL/plate-formes} :
\end{itemize}

\begin{verbatim}
  geospatial_lat_min      = geospatial_lat_max      = latitude  de la plateforme
  geospatial_lon_min      = geospatial_lon_max      = longitude de la plateforme
  geospatial_vertical_min = geospatial_vertical_max = altitude de la plateforme
\end{verbatim}

\begin{itemize}
\item
  pour une plateforme mobile :\\

  \begin{itemize}
  \item
    si l'option \texttt{coordinates} est fournie, alors les champs
    \textit{geospatial\_XXX} sont initialisés avec les données trouvées
    dans les fichiers Vecteurs des paramètres de la liste \texttt{coordinatess}.
  \item
    si l'option \texttt{coordinates} n'est pas fournie, \prodcdf
    affiche un warning.\\
    Les champs \textit{geospatial\_XXX} sont initialisés à \texttt{NaN}
    (ils doivent alors être initialisés manuellement dans le fichier de
    configuration).
  \end{itemize}
\end{itemize}


\paragraph{Champs \texttt{featureType}} : \label{ref:featureType}

Les spécificités de traitement liées à des options \texttt{featureType}
autre que \textit{Trajectory} n'ont pas (encore) été étudiées et
implémentées dans \prodcdf. (voir
\url{https://cfconventions.org/Data/cf-conventions/cf-conventions-1.10/cf-conventions.html\#discrete-sampling-geometries})

\subsubsection{Sections \textit{paramètre}}

\begin{center}
\tablefirsthead{}
\tablehead{\hline
  \textbf{option} &
  \textbf{nature} &
  \textbf{commentaire} \\}
\tabletail{\hline}
\tablelasttail{}
\begin{supertabular}{|p{3cm}|p{4cm}|p{9cm}|}
\hline
\textbf{option} &
\textbf{nature} &
\textbf{commentaire} \\
\hline

long\_name &
obligatoire &
pour une variable auxiliaire : \texttt{+modifier} ou \texttt{modifier+} \smallbreak
voir \nameref{ref:variables_auxiliaires}
\\ \hline

origin &
obligatoire &
nom du paramètre AIDA \smallbreak
(nom du vecteur \textbf{sans} son extension)
\\ \hline

ancillary\_variables &
obligatoire \smallbreak
(si des variables auxiliaires existent) &
liste des variables auxiliaires séparées par un espace
\\ \hline

standard\_name &
automatique &
pour une variable auxiliaire : \texttt{+modifier} ou \texttt{modifier+} \smallbreak
voir \nameref{ref:variables_auxiliaires} \smallbreak
voir \nameref{ref:note_standard_name}
\\ \hline

flag\_values \smallbreak
flag\_meanings &
obligatoire \smallbreak
(pour les variables auxiliaires de type \textbf{status\_flag}) &
voir \nameref{ref:variables_auxiliaires}
\\ \hline

units &
facultatif &
voir \nameref{ref:units}
\\ \hline

source &
facultatif &
voir \nameref{ref:source}
\\ \hline

\end{supertabular}
\end{center}

\begin{warning}
  les options qui ne sont pas explicitées ci-dessus sont inscrites dans la
  partie \textit{variables} du fichier \netcdf.
\end{warning}

\paragraph{Note \texttt{standard\_name}} : \label{ref:note_standard_name}

  Il peut arriver que la recherche automatique du \texttt{standard\_name} ne
  fonctionne pas. Dans ce cas, \prodcdf affiche un warning préfixé par
  \texttt{-+-}. Il est alors souhaitable de positionner l'option
  \texttt{standard\_name} dans le fichier de configuration.


\paragraph{Variables auxiliaires} : \label{ref:variables_auxiliaires}

Concernant l'utilisation et la description des variables auxiliaires
selon la norme CF, se reporter à \nameref{ref:annexes}.
\medbreak

Afin de faciliter la saisie du \texttt{long\_name} et du
\texttt{standard\_name} des variables auxiliaires dans le fichier de
configuration, ces 2 options peuvent être données sous les formes :

\begin{itemize}
\item
  \textbf{+\textit{modifier}}
\item
  \textbf{\textit{modifier}+}
\end{itemize}

A la génération du fichier \netcdf, le signe \texttt{+}
(plus) sera remplacé par le \texttt{long\_name} ou le
\texttt{standard\_name} de la \textit{variable principale}. Un espace sera
inséré entre les 2 parties du nom.
\medbreak

\textbf{Note} :

Le mot \textit{modifier} devrait être l'un de ceux décrits dans
la convention CF (voir \nameref{ref:annexes}).
\medbreak

\subparagraph{Exemple d'utilisation} : 

Si le fichier de configuration contient :

\begin{Verbatim}[fontsize=\myfontsize, commandchars=\\\{\}]
[SWD]
long_name           = Downward short wave radiation
origin              = ray_rgd_cnr1_c2_%60_Met_%1800
ancillary_variables = SWD_FLAG


long_name           = \textcolor{red}{+ancillary}
origin              = ctl_SB_invalid_cnr1_Ftp_%1800
standard_name       = \textcolor{blue}{+percentage_of_gap_filled_values}
\end{Verbatim}

sera retranscrit dans le fichier \netcdf sous la forme :

\begin{Verbatim}[fontsize=\myfontsize, commandchars=\\\{\}]
SWD:standard_name = "\textcolor{green}{surface_downwelling_shortwave_flux_in_air}" ;

SWD_FLAG:long_name = "Downward short wave radiation \textcolor{red}{ancillary}" ;
SWD_FLAG:standard_name = "\textcolor{green}{surface_downwelling_shortwave_flux_in_air} \textcolor{blue}{percentage_of_gap_filled_values}" ;
\end{Verbatim}

\begin{warning}
Attention ! \\
Les variables auxiliaires \textbf{ne sont pas toutes de type
\textit{status\_flag}}. Ce type de variable auxiliaire est réservé à des
variables qui ont un nombre fini de valeurs possibles. Leur descritpion
doit donc contenir les 2 options \texttt{flag\_values} et
\texttt{flag\_meanings} (voir \nameref{ref:annexes}).
\end{warning}


\paragraph{Unité d'un paramètre} : \label{ref:units}

L'unité des paramètres est lue dans le fichier Vecteur associé.
Cependant elle peut être positionnée par la configuration de l'option
\texttt{units}.\\
Si l'unité lue ou configurée est \texttt{unitless} (ou pour des raisons
historiques \texttt{unknown} ou \texttt{without}), elle est inscrite dans le
fichier \netcdf avec la valeur : \texttt{1} (standard UDUNITS).


\paragraph{Source d'un paramètre} : \label{ref:source}

Les capteurs sources d'un paramètre sont recherchés automatiquement par
\prodcdf par interrogation l'outil AIDA
\textbf{\texttt{gengraphvect.py}}. Cependant, elle peut être positionnée par la
configuration de l'option \texttt{source} (en particulier si
\texttt{gengraphvect.py} n'arrive pas à touver l'information, ce qui
peut se produire si un vecteur a été copié depuis une autre
acquisition).
\medbreak

Le format préconisé pour l'option \texttt{source} est :\\
\begin{Verbatim}[fontsize=\myfontsize]
fabricant_1 : modèle_1 : numéro_de_série_1; fabricant_2 : modèle_2 : numéro_de_série_2
\end{Verbatim}

Ex :\\

\begin{Verbatim}[fontsize=\myfontsize]
source = Atexis : PT1000 : 0001; Vaisala : HMP110 : L4150552; Vaisala : PTB220 : X2050007
\end{Verbatim}

\subsubsection{Options libres}\label{ref:options_libres}

  Les \textit{options libres} positionnées dans la section
  \texttt{[GENERAL]} du fichier de configuration seront
  retranscrites sans vérification, dans la partie \textit{global
  attributes} du fichier \netcdf.
  \medbreak

  Les \textit{options libres} positionnées dans les sections \textit{paramètre}
  du fichier de configuration seront retranscrites sans vérification, dans
  la partie \textit{variables} du fichier \netcdf.


\paragraph{Format} :

Les \textit{options libres} sont retranscrites sous forme de
\textbf{chaîne de caractères} dans le fichier \netcdf.

Il est cependant possible, pour \textbf{les sections \textit{paramètre}
seulement}, de forcer l'inscription \textbf{au même format} que celui
du vecteur de données. Pour cela, il suffit d'ajouter le mot-clé \texttt{:NUM}
au nom de l'option.
\medbreak

Exemples :

\begin{longtable}[]{@{}ll@{}}
\toprule
\begin{minipage}[b]{0.47\columnwidth}\raggedright
fichier de configuration\strut
\end{minipage} & \begin{minipage}[b]{0.47\columnwidth}\raggedright
fichier \netcdf\strut
\end{minipage}\tabularnewline
\midrule
\endhead
\begin{minipage}[t]{0.47\columnwidth}\raggedright
valid\_min = 0.\strut
\end{minipage} & \begin{minipage}[t]{0.47\columnwidth}\raggedright
valid\_min = "0."\strut
\end{minipage}\tabularnewline
\begin{minipage}[t]{0.47\columnwidth}\raggedright
valid\_min:NUM = 0\strut
\end{minipage} & \begin{minipage}[t]{0.47\columnwidth}\raggedright
valid\_min = 0.f\strut
\end{minipage}\tabularnewline
\begin{minipage}[t]{0.47\columnwidth}\raggedright
valid\_range:NUM = 0., 1000.\strut
\end{minipage} & \begin{minipage}[t]{0.47\columnwidth}\raggedright
valid\_range = 0.f, 1000.f\strut
\end{minipage}\tabularnewline
\bottomrule
\end{longtable}
(les affichages pour la colonne 'fichier \netcdf' sont ceux donnés par
l'outil \texttt{ncdump})
\medbreak

Il est également possible de forcer une inscription \textbf{dans un
autre format numérique} que celui du vecteur de données (bien que ceci
pourrait ne pas avoir beaucoup de sens\ldots\hspace{0pt}). Pour cela, il
suffit d'ajouter \texttt{:fmt} au nom de l'option, avec pour valeur de
\texttt{fmt} une valeur parmi : \texttt{int8}, \texttt{int16},
\texttt{int32}, \texttt{int64}, \texttt{float32}, \texttt{float64}.

Exemples :

\begin{longtable}[]{@{}ll@{}}
\toprule
\begin{minipage}[b]{0.47\columnwidth}\raggedright
fichier de configuration\strut
\end{minipage} & \begin{minipage}[b]{0.47\columnwidth}\raggedright
fichier \netcdf\strut
\end{minipage}\tabularnewline
\midrule
\endhead
\begin{minipage}[t]{0.47\columnwidth}\raggedright
valid\_range:int8 = 0, 2\strut
\end{minipage} & \begin{minipage}[t]{0.47\columnwidth}\raggedright
valid\_range = 0b, 2b ;\strut
\end{minipage}\tabularnewline
\begin{minipage}[t]{0.47\columnwidth}\raggedright
flag\_values:int8 = 0, 1, 2\strut
\end{minipage} & \begin{minipage}[t]{0.47\columnwidth}\raggedright
flag\_values = 0b, 1b, 2b ;\strut
\end{minipage}\tabularnewline
\begin{minipage}[t]{0.47\columnwidth}\raggedright
flag\_meanings = bon, moyen, mauvais\strut
\end{minipage} & \begin{minipage}[t]{0.47\columnwidth}\raggedright
flag\_meanings = "bon, moyen, mauvais" ;\strut
\end{minipage}\tabularnewline
\bottomrule
\end{longtable}
(les affichages pour la colonne 'fichier \netcdf' sont ceux donnés par
l'outil \texttt{ncdump})
\medbreak


\subsection{Sélection de vecteur}\label{ref:vector_select}

La mécanique de \textit{sélection de vecteur} permet de sélectionner le
vecteur (et éventuellement son commentaire assissié) à utiliser pour un
paramètre \textbf{en fonction d'un acqid}.\\
\medbreak

Exemple : utilisation du vecteur d'altitude généré par l'instrument
\textit{aipov} ou par l'instrument \textit{gps} en fonction de l'acqid
traité.
\medbreak

La mise en oeuvre de la mécanique de \textit{sélection de vecteur} se fait
dans un fichier de type INI (voir \nameref{sec:anx_ini}).

\subsubsection{Formalisme}

Chaque section décrit un choix possible.  \\

Chaque section \textbf{doit obligatoirement} contenir un couple d'options
\textbf{\textit{default}} et \textbf{\textit{default\_cmt}} : ces
options seront utilisées s'il n'existe pas de ligne spécifique pour
l'acqid traité. \\

Chaque section peut contenir plusieurs couples d'options :
\begin{itemize}
  \item \textbf{\textit{XXX}} = valeur
  \item \textbf{\textit{XXX\_cmt}} = valeur
\end{itemize}
\medbreak

où \textbf{\textit{XXX}} peut prendre l'une des trois formes suivantes :

\begin{tabular}{|l|p{8cm}|}
\hline
 acqid & un acqid unique
 
 (ex: Tf230001) \\
\hline
 acqid1-acqid2 & une liste d'acqids compris entre acqid1 et acqid2
 
 (ex Tf230001-Tf23005) \\
\hline
 acqid1-* & une liste d'acqids compris entre acqid1 et le dernier acqid à
 traiter
 
 (ex Tf230001-*) \\
\hline
\end{tabular}
\medbreak

Ces options (\textbf{\textit{XXX}} et \textbf{\textit{XXX\_cmt}}) seront
utilisées pour sélectionner les valeurs spécifiques pour
l'acqid ou la liste d'acqids traité(s).
\medbreak

\subsubsection{Utilisation dans le fichier de configuration de \prodcdf}

La syntaxe d'utilisation dans le fichier de configuation de
\prodcdf est la suivante : \\
- pour sélectioner la variable :
  \texttt{\textbf{\textcolor{red}{\$\{section\}}}} \\
- pour sélectioner le commentaire de la variable :
  \texttt{\textbf{\textcolor{red}{\$\{section\_cmt}\}}}
\medbreak

\subsubsection{Exemple}

Fichier de sélection :
\begin{Verbatim}[fontsize=\myfontsize, commandchars=\\\{\}]
[choix_lat]
default      : gps
default_cmt  : le GPS est plus précis dans les sondages que l'AIPOV
fs170012     : aipov
fs170012_cmt : utilisation de l'AIPOV car le GPS est HS
\end{Verbatim}

Fichier de configuration de \prodcdf :
\begin{Verbatim}[fontsize=\myfontsize, commandchars=\\\{\}]
[LAT]
long_name = latitude $\{choix_lat\}
origin    = pos_lat_$\{choix_lat\}_1
comment   = $\{choix_lat_cmt\}
\end{Verbatim}


\subsection{Exemples de fichiers de configuration de \prodcdf}

\newgeometry{margin=1cm}
\begin{landscape}

\subsubsection{Mode \texttt{periode}}

Excécution :

\begin{Verbatim}[fontsize=\myfontsize, commandchars=\\\{\}]
prod_CDF.py -c tests/Tf-AERIS-base.V7.ccdf -S -o /tmp/Tf -v 1 Tf210001 Tf210365
\end{Verbatim}

Extrait du fichier Tf-AERIS-base.V7.ccdf

\begin{Verbatim}[fontsize=\myfontsize, commandchars=\\\{\}]
[DIRECTIVES]
prefix                 = CNRM_MTO_L2A-30min-Mean
plateforme             = toulouse-flux
reftime                = rft_time_REG_%1800
mode                   = periode
date_format            = %Y-%m-%d
replace_missing_vector = True

[GENERAL]
title        = Meteopole-Flux Meteorological Data
summary      = This dataset contains air temperature, humidity at 2 levels...
keywords     = earth science, atmosphere, atmospheric radiation, ...
doi          = https://doi.org/10.25326/44
source       = ground station

creator_name = X Y

contributor_name = X1: Y1: X1.Y1@domain.fr: Institute;
                   X2: Y2: X2.Y2@domain.fr: Institute;
                   X3: Y3: X3.Y3@domain.fr: Institute;

contributor_role = principalinvestigator;
                   resourceprovider;
                   contributor;

...

[SWD]
long_name           = Downward short wave radiation
origin              = ray_rgd_cnr1_c2_%60_Met_%1800
ancillary_variables = SWD_FLAG

[SWD_FLAG]
long_name           = +ancillary
origin              = ctl_flag_invalid_cnr1_Ftp_%1800
standard_name       = +percentage_of_gap_filled_values
...

# Pourcentage de points de visi sur 30 min < 60km OK
[VISI_MEAS_RATE]
long_name           = Rate of visibility measurement in valid measurement range
origin              = opt_pom_dso_tr30_%60_per_%1800
comment             = Rate of measurements in valid (0, 59999 m), measurement range computed during 30 min

# Moyenne harmonique sur 30mn
[VISI_MOYH]
long_name           = Visibility's harmonic mean
origin              = opt_pom_dso_tr30_%60_Moyh_%1800
valid_range:NUM     = 0, 59999
comment             = Harmonic mean of visibility computed during 30 minutes with measurements into valid range (0,  59999 m). If one of the values was out of valid range the result is set to -99
\end{Verbatim}

\subsubsection{Mode \texttt{mobile}}

Excécution :

\begin{Verbatim}[fontsize=\myfontsize, commandchars=\\\{\}]
prod_CDF.py -c tests/prod_g2401m_%60.ccdf -S tests/choix_vecteurs_params -o /tmp/as -v 1 as220024 as220025
\end{Verbatim}

Extrait du fichier tests/prod\_g2401m\_\%60.ccdf

\begin{Verbatim}[fontsize=\myfontsize, commandchars=\\\{\}]
[DIRECTIVES]
prefix     = ACROSSS-2022_SAFIRE-ATR42_SAFIRE_CHEMISTRY-PICARRO_1MN
plateforme = atr
reftime    = rft_time_PICARRO_%60
mode       = mobile


[GENERAL]
title        = ACROSS-2022 data from Picarro instrument G2401-m
summary      = SAFIRE ATR42 Picarro-G2401-m in situ measurements
keywords     = earth science, atmosphere, aircraft, aircraft motion sensor,
               altitude, CH4, CO2, CO
project      = ACROSS-2022
doi          = http://dx.doi.org/10.25326/467

source       = airborne observation
platform     = SAFIRE-ATR42
coordinates  = time LONGITUDE LATITUDE ALTITUDE

creator_name = X1 Y1, X2, Y2
comment      = Warning : Most measurements are not valid before take-off and
               after landing
history      = First version of Picarro-g2401-m data

contributor_name = X3:Y3:X3.Y3@domain.fr;
                   X4:Y4:X4.Y4@domain.fr;
contributor_role = resourceprovider;
                   principalinvestigator;


[LATITUDE]
long_name           = Latitude
#origin              = pos_lat_g2401m_valid_cal_%60
origin              = ${choix_lat}

[LONGITUDE]
long_name           = Longitude
origin              = pos_lon_g2401m_valid_cal_%60

[ALTITUDE]
long_name           = Altitude
origin              = alt_g2401m_valid_cal_%60
standard_name       = altitude

[CH4]
long_name           = CH4 mixing ratio from Picarro instrument G2401-m
origin              = chm_CH4_g2401m_valid_cal_%60
ancillary_variables = CH4_STDEV FLAG_CH4 NB_POINT_CH4

[CH4_STDEV]
long_name           = Standard deviation of CH4 mixing ratio from Picarro instrument G2401-m
origin              = chm_CH4_stdev_g2401m_valid_cal_%60

[NB_POINT_CH4]
long_name           = Number of points used for the calculation CH4 mixing ratio from Picarro instrument G2401-m
origin              = chm_CH4_NbPoints_g2401m_valid_cal_%60

[FLAG_CH4]
long_name           = Flag+
origin              = chm_CH4_flag_g2401m_valid_cal_%60
standard_name       = +status_flag
flag_values         = U, N, O, K, R, H
flag_meanings       = data_correct_before_manual_quality_control data_incorrect_before_manual_quality_control data_correct_after_manual_quality_control data_incorrect_after_manual_quality_control data_correct_after_manual_quality_control_and_backwards_propagation data_incorrect_after_manual_quality_control_and_backwards_propagation

...
\end{Verbatim}


Fichier de sélection passé en argument \texttt{choix\_vecteurs\_params}

\begin{Verbatim}[fontsize=\myfontsize, commandchars=\\\{\}]
# utilisation dans le fichier de configuration de prod_CDF.py :
#   origine = ven_e_${choix_lat}_25
[choix_lat]
default      = pos_lat_g2401m_valid_cal_%60
default_cmt  = test de prod_CDF
fs170012     = aipov
fs170012_cmt = utilisation de l'AIPOV car le GPS est HS
as220025     = pos_lat_XXX_valid_cal_%60
as220025_cmt = validation de XXX
\end{Verbatim}

\subsubsection{Type \texttt{status\_flag}}

Extrait du fichier as\_status\_flag.V1.ccdf

\begin{Verbatim}[fontsize=\myfontsize, commandchars=\\\{\}]
[DIRECTIVES]
prefix     = as_status_flag
plateforme = atr
reftime    = rft_time_PICARRO_%60
mode       = mobile

...

[CH4]
long_name           = CH4 mixing ratio from Picarro instrument G2401-m
origin              = chm_g2401m_valid_cal_%60
ancillary_variables = FLAG_CH4

[FLAG_CH4]
long_name           = Flag+
origin              = chm_flag_g2401m_valid_cal_%60
standard_name       = +status_flag
flag_values         = U, N, O, K, R, H
flag_meanings       = data_correct_before_manual_quality_control data_incorrect_before_manual_quality_control data_correct_after_manual_quality_control data_incorrect_after_manual_quality_control data_correct_after_manual_quality_control_and_backwards_propagation data_incorrect_after_manual_quality_control_and_backwards_propagation
\end{Verbatim}

\end{landscape}
\restoregeometry


\subsection{Annexes NetCDF}\label{ref:annexes}

\subsubsection{Variables auxiliaires}


\paragraph{\texttt{long\_name} et \texttt{standard\_name}} :

Voir :
\url{https://cfconventions.org/Data/cf-conventions/cf-conventions-1.10/cf-conventions.html\#ancillary-data}

Le \texttt{long\_name} et le \texttt{standard\_name} d'une variable
auxiliaire devraient être ceux de sa \textit{variable principale} augmenté
d'un \textbf{modificateur}

Exemple :

\begin{itemize}
\item
  pour la variable \texttt{q}\\
  standard\_name = \textbf{specific\_humidity}\\
  ancillary\_variables = \texttt{q\_error\_limit}
  \texttt{q\_detection\_limit}
\item
  pour la variable auxiliaire \texttt{q\_error\_limit}\\
  standard\_name = \textbf{specific\_humidity} {standard\_error}
\item
  pour la variable auxiliaire \texttt{q\_detection\_limit}\\
  standard\_name = \textbf{specific\_humidity} {detection\_minimum}
\end{itemize}

La description des 4 modificateurs prévus par la convention CF est
présisée sous :\\
\url{https://cfconventions.org/Data/cf-conventions/cf-conventions-1.10/cf-conventions.html\#standard-name-modifiers}.\\
Ces modificateurs sont :

\begin{itemize}
\item
  \textit{detection\_minimum}
\item
  \textit{number\_of\_observations}
\item
  \textit{standard\_error}
\item
  \textit{status\_flag}
\end{itemize}

Le modificateur \textit{status\_flag} est un cas particulier : voir
\nameref{ref:status_flag}


\subparagraph{Attention} :

Cette liste ne couvre pas tout les besoins (ex :
variable auxiliaire précisant le pourcentage de valeurs remplacées). De ce
fait, \prodcdf ne vérifie pas les modificateurs qui lui sont fournis
et \textbf{l'utilisateur peut définir son propre modicateur} sous la
forme de mots séparés par le caractère '\_' (ex :
percentage\_of\_gap\_filled\_values)


\paragraph{\texttt{status\_flag}} : \label{ref:status_flag}

Voir :\\
\url{https://cfconventions.org/Data/cf-conventions/cf-conventions-1.10/cf-conventions.html\#flags}~\\
et\\
\url{https://cfconventions.org/Data/cf-conventions/cf-conventions-1.10/cf-conventions.html\#standard-name-modifiers}

Toutes les variables auxiliaires ne sont pas de type
\textit{status\_flag}.\\
Dans le cas où la variable auxiliaire est vraiment un
\textit{status\_flag}, la section de description de la variable doit
contenir les options \texttt{flag\_values} et \texttt{flag\_meanings} :

\begin{itemize}
\item
  \texttt{flag\_values} est l'ensemble des valeurs possibles de la
  variable.\\
  Les valeurs possibles sont séparées par des virgules.
\item
  \texttt{flag\_meanings} est une description de chacune des valeurs de
  \texttt{flag\_values}.\\
  La descritpion d'une valeur est un ensemble de mots séparés par des
  '\_'.\\
  Les descriptionis sont séparées par des espaces.
\end{itemize}

Exemples :

\begin{verbatim}
flag_values:int8 = 0, 1, 2
flag_meanings    = bon moyen mauvais
\end{verbatim}

\begin{verbatim}
flag_values   = U, N, O, K, R, H
flag_meanings = data_correct_before_manual_quality_control data_incorrect[...]
\end{verbatim}

%%:::::::::: ICI début
%%:::::::::: ICI fin
\begin{warning}

  Pour les variables auxiliaires de type \textit{status\_flag} pour
  lesquelles les valeurs possibles (\textit{flag\_values}) sont des lettres
  (cf. lexemple ci-dessus), l'option \textit{flag\_values} est enregistrée
  par \prodcdf dans le fichier \netcdf sous la forme Python "liste de lettres".
  Cependant certains outils de validation interprètent ce format comme
  "une chaîne de caractères" et peuvent afficher une erreur.
\end{warning}

\subsubsection{Validation des fichiers \netcdf}

Il existe plusieurs outils pour valider les fichiers
\netcdf dont : \url{https://compliance.ioos.us/index.html}

\subsubsection{Visualisation des fichiers \netcdf}

Il existe plusieurs outils de visualisation des fichiers
\netcdf dont :

\begin{itemize}
\item
  ncdump : \textit{dump} d'un fichier \netcdf en texte.
\item
  ncview et panoply : outils graphiques
\end{itemize}

